\documentclass[11pt,a4paper]{article}
\usepackage[T1]{fontenc}
\usepackage[utf8]{inputenc}
\usepackage{amsmath,amssymb,xcolor}
\usepackage[margin=25mm]{geometry}
\usepackage[colorlinks=true,urlcolor=blue]{hyperref}
\definecolor{nexorablue}{HTML}{2459B5}
\setlength{\parindent}{0pt}
\setlength{\parskip}{7pt}
\setlength{\emergencystretch}{2em}
\renewcommand{\familydefault}{\sfdefault}
\begin{document}

{\small\bfseries\color{nexorablue} NEXORA CPGE}\par
{\small MPSI / KUT-SUP-06-01 / Énoncé}\par
{\LARGE\bfseries Trois quotients complexes}\par
\vspace{0.25em}\hrule\vspace{0.6em}

Écrire les trois quotients sous la forme $a+ib$, avec $a,b\in\mathbb R$.

\[ \frac{1}{i},\qquad \frac{2-i}{3+4i},\qquad \frac{1-2i}{-2+5i}. \]

\bigskip
{\footnotesize Adaptation française par Nexora CPGE. Source : Ken Kuttler (Lyryx),
\href{https://math.libretexts.org/Bookshelves/Linear_Algebra/A_First_Course_in_Linear_Algebra_(Kuttler)/06\%253A_Complex_Numbers/6.01\%253A_Complex_Numbers}{\textit{A First Course in Linear Algebra}}, § 6.1, exemple 3.
\href{https://creativecommons.org/licenses/by/4.0/}{CC BY 4.0}.}
\end{document}
